\documentclass[a4paper, 12pt]{article}
\usepackage{ulem}
\parindent 0px
\usepackage{amsmath, amsfonts, amssymb}
\usepackage{float}
\usepackage[margin=1cm]{geometry}
\title{MATRIX}
\author{PHILIP THEURI}
\date{9th June 2026}
\begin{document}
	\maketitle
	\section*{\huge\centering\uuline{ MATRIX}}
	\begin{equation*}
A=
\begin{matrix}
	a & b\\
	c & d\\
	e & f\\
	\end{matrix}
	\end{equation*}
\begin{equation*}
	A=	
	\begin{pmatrix}
a & b & c\\
d  & e & f \\
\end{pmatrix}	
\end{equation*}	
	
	\begin{equation*}
		B=
	\begin{bmatrix}
		1 & 2 & 3 & 4 \\
		5 & 6 & 0 & 7\\
	\end{bmatrix}
	\end{equation*}
	
	
		\begin{equation*}
		B=
		\begin{Bmatrix}
			1 & 2 & 3 & 4 \\
			5 & 6 & 0 & 7\\
		\end{Bmatrix}
	\end{equation*}
	
	
	\begin{equation*}
		C=
	\begin{vmatrix}
		2 & 3 & b & c & v & 3\\
		8 & 7 & 4 & 2 & 0 & 6\\
		Z & d & b & 9 & h & 2\\
		1 & 7 & 9 & 9 & h & o\\
		\end{vmatrix}
	\end{equation*}
	
	
		\begin{equation*}
		C=
		\begin{Vmatrix}
			2 & 3 & b & c & v & 3\\
			8 & 7 & 4 & 2 & 0 & 6\\
			Z & d & b & 9 & h & 2\\
			1 & 7 & 9 & 9 & h & o\\
		\end{Vmatrix}
	\end{equation*}
	
	
	\begin{equation}
		D=
		\begin{vmatrix}
		a_{11} &  a_{12}& a_{13}\\
		a_{21} & a_{22} & a_{23}\\
		a_{31} & a_{32} & a_{33}\\
		\end{vmatrix}
		\end{equation}
	for addition or subtraction of two or more 2 just write the 2 equations in between the the opening and closing equation.
	\begin{equation*}
		D=
		\begin{vmatrix}
			a_{11} &  a_{12}& a_{13}\\
			a_{21} & a_{22} & a_{23}\\
			a_{31} & a_{32} & a_{33}
		\end{vmatrix}
+
		\begin{Vmatrix}
			2 & 3 & b & c & v & 3\\
			8 & 7 & 4 & 2 & 0 & 6\\
			Z & d & b & 9 & h & 2\\
			1 & 7 & 9 & 9 & h & o
		\end{Vmatrix}
	\end{equation*}
	
	
	
	\begin{gather*}
	2x+y =20\\
	3x -z=20
	\end{gather*}
	
	
		
	\begin{align*}
		2x+y &=20\\
		3x -z&=20
	\end{align*}
	
	\begin{align*}
		2x&+y &=20&\\
		3x& - z &=20&
	\end{align*}
	\pagebreak
	\section*{\huge\centering\uuline{WRITING EQUATIONS}}
	There are different ways of writing an equation;
	
	\begin{itemize}
\item\begin{equation*}
	y=mx + C_1;
	\end{equation*}
\item\begin{equation*}
	y= x^2 + \sqrt{x};
	\end{equation*}
	\item	$$y= x^2 + \sqrt{x};$$
		\item 	\[y= x^2 + \sqrt{x};\]\\
			\end{itemize}
			
		$$y=\sum_{n=0}^{N} \mathrm{n^2} $$	
		$$y = \sum_{n} \mathrm{x^{2}_{n}};$$
		
		$$y =\int xdx$$
			$$y = \int\limits_{-\infty}^{\infty} x\mathrm{dx}$$
			
		$$y =\int xdx$$
		$$y = \int_{-\infty}^{\infty} x\mathrm{dx}$$	
		\subsection*{For a piece-wise equation;}
	\Large{	\[F_x=
		\begin{cases}
			o, & \text{if} \,\, x < a;\\
			\frac{x-a}{b-a}, & \text{if} \,\, a \,\leq  b<a\\
			1, & \text{if}\,\, x\geq b
			\end{cases}
		 \]	
	}	
	$$\phi$$
	\[\hat{y}=\tilde{x}+ \bar{z}\]
	
	\[\gamma= \alpha + \beta + \pi +\psi \]
\end{document}